\documentclass[11pt]{article}
\usepackage[margin=1in]{geometry}
\usepackage{times}
\usepackage{hyperref}
\hypersetup{colorlinks=true, linkcolor=black, urlcolor=blue, citecolor=black}
\title{Spin Without a Source: A Critique of Torsion Field Claims as a Gravitational Mechanism}
\author{Flyxion \\ \small Independent Researcher}
\date{}
\begin{document}
\maketitle

\begin{abstract}
Torsion field theories, most prominently associated with the work of Anatoly Akimov and Gennady Shipov, propose a subtle, information-carrying field generated by rotating or spinning matter, distinct from gravitation and electromagnetism, capable of instantaneous transmission and, in some formulations, of locally altering weight or inertia. This paper distinguishes the legitimate mathematical object of torsion in differential geometry, a genuine and well-understood generalization of Riemannian curvature used in some extensions of general relativity, from the informal, empirically unconstrained "torsion field" of the Akimov-Shipov tradition, and argues that the latter borrows the former's name and superficial vocabulary while dispensing with everything that made it mathematically or observationally accountable.
\end{abstract}

\section{Introduction}

Torsion, in the technical sense used in Einstein-Cartan theory and related extensions of general relativity, is the antisymmetric part of a connection's coefficients, a well-defined geometric object that couples to the intrinsic spin density of matter rather than to its mass-energy alone, producing effects that are exceedingly small under ordinary conditions and have never been detected, consistent with theoretical expectation, at any laboratory or astrophysical scale probed so far. This is a legitimate and actively studied area of theoretical physics, if a comparatively minor one, with a small effect size predicted precisely because the coupling constant governing it is tied to the same order of magnitude as other quantum gravitational corrections.

The torsion field of Akimov and Shipov's popular and commercial work, sometimes called spin field or axion field in loosely translated Russian-language literature from the 1990s, shares essentially nothing with this except the word torsion. It is described as generated by any rotating object, propagating without attenuation or delay across arbitrary distances, carrying information rather than energy, and capable of influencing biological systems, consciousness, and the weight of nearby objects. None of these properties follows from, or is consistent with, the differential-geometric object the name is borrowed from.

\section{Why "Carries Information Without Carrying Energy" Is Not a Loophole}

A recurring defense of the informal torsion field is that its undetectability by conventional instruments is expected, since it is claimed to carry information rather than energy, and instruments built to detect energy transfer would naturally miss a signal of this kind. This framing treats the absence of energy transfer as a feature that explains away the absence of detection, but information transfer without any accompanying physical interaction is not a subtle experimental challenge, it is a description of something that cannot interact with any detector at all, by definition, since detection is itself a physical interaction. A field that in principle cannot be measured by any physical process is not describing a hidden layer of physics; it is describing the absence of a testable claim.

\section{The Entropic Gravity Framing}

On an account where gravitational curvature tracks the redistribution of an ordered configuration across coupled scalar, vector, and entropy fields, weight and inertia are not independently adjustable quantities that a separate, parallel field could tune from outside; they are expressions of how a given mass-energy distribution participates in that redistribution. A device claiming to alter local weight through spin-generated torsion fields would need a mechanism by which ordinary mechanical rotation, itself an entirely mundane redistribution of angular momentum with no unusual entropy signature, somehow taps into and modifies the much larger-scale bookkeeping that produces gravitational curvature. No such channel exists in this framework, any more than in standard general relativity, and the informal torsion field literature has never proposed one with enough specificity to evaluate, relying instead on qualitative language about subtle energies and information fields that resists translation into any testable prediction.

\section{Objections}

Proponents sometimes cite anecdotal or small-sample biological and materials-science experiments reporting effects from torsion generators, including changes in crystal structure, seed germination rates, or reaction times in human subjects. These studies, where they exist in published form, are characterized by small sample sizes, absence of blinding, and no independent replication by laboratories outside the group that produced the original device, a pattern indistinguishable from publication bias and experimenter effect rather than a genuine physical signal.

\section{Conclusion}

The differential-geometric object called torsion is real, well defined, and predicted to produce effects far too small to explain any of the claims made under its borrowed name. The Akimov-Shipov torsion field, evaluated on its own terms, describes an interaction with no available mechanism to alter gravitational curvature and no falsifiable prediction that has survived independent scrutiny, and the informational framing offered in its defense removes the claim from the category of physics rather than protecting it within it.

\end{document}
